% =============================================================================
%  Sección 2: Desarrollo
%  Sustituye el contenido de \lipsum con tu propio texto.
%  Agrega o elimina subsecciones según la estructura de tu argumento.
% =============================================================================

\section{Ru- Urbanismo (Gobierno)}

\subsection{El PRI antes de la Coyuntura}
\label{PRI-pre-Coyuntura}

Entender el desarrollo de un país tan diverso como lo es México es imposible sin analizar su
historia moderna y contemporánea. A diferencia de otros países del continente latinoamericano, 
los Estados Unidos Mexicanos solo atravesaron un movimiento social que escaló a ser una guerra civil;
La Revolución Mexicana. La cual dió paso a la instauración de una dictadura única en el continente. No 
era dirigida por un hombre, si no por un partido político con tintes de izquierda, derecha y centro 
a conveniencia. Sin embargo, durante buena parte del siglo XX, el país avanzó de manera asombrosa, logró 
recuperarse a niveles históricos y sus PIB nacional creció a un promedio de un 2\,\% anual durante los años 
1940 a 1970. Fue durante este periodo que nació la clase media mexicana, mayormente dominada por trabajadores, 
pequeños burgueses y uno que otro agricultor. Y con ello las urbes mexicanas tuvieron un \termino{boom} en 
términos geográficos y económicos. Las ciudades pasaron ser pequeños nodos aislados a ser verdaderas metropolis 
interconectadas, dadoras de bienes y servicios a la naciente clase media obrera mexicana. El actor detrás de este 
sistema fue el mismo gobierno, instaurando así lo que se conocería como \termino{Estado de Bienestar}. El estado 
mexicano se convirtió en el dador de servicios públicos (Salud, vivienda, electricidad, agua corriente, entretenimiento, etc.) 
y para todo ese sistema se necesitaba mano de obra; convirtiendo así al partido en el poder en un negocio redondo. 
La maquinaria política afinó sus engranajes y el partido hegemónico absorvió el sistema económico, tecnológico, 
académico, sindical, militar, de entretenimiento y posteriormente la prensa y la expresión de sus ciudadanos. \cite{agustin_tragicomedia1} Sin embargo, 
este logro no fue único de México; el sistema político mundial viró hacia esa política: Se necesitaba de un 
un régimen que cumpliera con mantener un órden político y social aunque fuera represivo y autoritario. 
Los componentes para ello no fueron implementados por ocurrencias, por lo menos en México durante la primera mitad del 
siglo XX se temía un movimiento armado a escala de la Revolución Mexicana acontecida apenas un par de décadas atrás. 
Por lo que mantener el sistema requirió de actores especiales: 
\begin{itemize}
    \item \textbf{L{\'a}zaro C{\'a}rdenas Del R{\'i}o:} El caudillo que se convirti{\'o} en presidente y desplaz{\'o} al antiguo r{\'e}gimen de 
          Plutarco El{\'i}as Calles, instaurando un nuevo \textbf{ordenamiento pol{\'i}tico} con una labor noble excepcional: la reivindicaci{\'o}n 
          de la lucha social de las clases m{\'a}s desprotegidas, como lo eran los campesinos y obreros.
    
    \item \textbf{Jos{\'e} Vasconcelos:} Impulsor del \textbf{aparato educativo y cultural} posrevolucionario, quien mediante la fundaci{\'o}n de la SEP y 
          las misiones culturales busc{\'o} integrar a las masas rurales a la identidad nacional.
    
    \item \textbf{Carlos Contreras:} Pionero de la planificaci{\'o}n urbana moderna en M{\'e}xico, responsable de proyectar la 
          \textbf{configuraci{\'o}n estructural} y vial de la capital para adaptarla a la incipiente industrializaci{\'o}n del siglo XX.
    
    \item \textbf{Antonio Ortiz Mena:} Arquitecto financiero del ``Milagro Mexicano'', cuyo \textbf{esquema econ{\'o}mico} de Desarrollo 
          Estabilizador logr{\'o} un crecimiento sostenido del PIB con estabilidad cambiaria y baja inflaci{\'o}n.
    
    \item \textbf{Fidel Vel{\'a}zquez:} L{\'i}der corporativista que consolid{\'o} el \textbf{entramado sindical} de la CTM, convirti{\'o} 
          al movimiento obrero en un pilar de estabilidad para el r{\'e}gimen y negoci{\'o} el acceso a prestaciones sociales institucionales.
    
    \item \textbf{Bernardo Quintana Arrioja:} Ingeniero civil y empresario clave en la creaci{\'o}n del \textbf{complejo de infraestructura} contempor{\'a}neo 
          del pa{\'i}s~\cite{fundacionica2023viaje}.
    \begin{itemize}
        \item \textbf{Fundaci{\'o}n de ICA (1947):} A los 28 a{\~n}os reuni{\'o} a 17 ingenieros de la UNAM para constituir Ingenieros Civiles Asociados, 
              transformando la obra p{\'u}blica mexicana mediante una escuela de alta calidad t{\'e}cnica y gerencial~\cite{fundacionica2023viaje}.
        \item \textbf{Implementaci{\'o}n de la Soluci{\'o}n Mil{\'a}n:} Introdujo los muros de contenci{\'o}n ``Mil{\'a}n'' 
              durante la planeaci{\'o}n y construcci{\'o}n de las primeras l{\'i}neas del Metro de la CDMX. Este hito resolvi{\'o} el reto estructural de 
              excavar t{\'u}neles bajo las desfavorables condiciones del suelo arcilloso del subsuelo lacustre de la capital, evitando colapsos urbanos~\cite{fundacionica2023viaje}.
    \end{itemize}

\end{itemize}

No se puede negar que el regimen político de aquellos años contenía los problemas que sufrimos hoy día (nepotismo, corrupción, escándalos políticos y artísticos 
y un largo etc.), sin embargo, la dictadura de un solo partido logró convertir varios obstáculos y necesidades ciudadanas en poderosas secretarías, 
a la par que brindaba salud, movilidad, vivienda y demás bienes y servicios a la población de las metropolis circundantes. A pesar de que estos bienes 
solo estaban disponibles para quienes tuvieran el privilegio de venir a la capital y establecerse en alguna de las alcaldías centrales (Cuahtémoc, Benito Juárez)
o el Norte de la Ciudad de México. Para el resto del país, la situación seguiría mas o menos igual. \\
Aquí es un buen momento para detenernos a pensar y analizar con lupa \textit{¿Acaso el viejo PRI representó el avance más significativo del país?}, la respuesta 
es un rotundo \textit{NO} con sus respectivos negro-grises. El estado de bienestar, instaurado después de la Segunda Guerra Mundial y heredado por los 
Estados Unidos de América no fue una política exclusiva del partido hegemónico de esos años, fue una política heredada por todo el bloque occidental y por algunos 
países dentro de la unión soviética (de donde tomó su origen). Dicho sistema se basaba en un estado único, fuerte y poderoso que debía administrar y controlar los 
tres pilares de un país: Sociedad, política y economía \cite{bauman_estado}; el PRI destacó por crear secretarías, burocracía y controlar los movimientos obreros a manos de los 
\termino{Tres Lobitos} y su Lobito principal: Fidel Velázquez. Quién fungió por más de 40 años como el líder sindical más longevo de la historia de México \cite{agustin_tragicomedia1}.
Pero este no sirvió a los intereses de los trabajadores o al sindicato como ente protector de la fuerza obrera mexicana; todo lo contrario, sirvió a los intereses del estado 
hegemónico, ganadose el apodo de \termino{Líder Charro}. \\
Existe una disconcordancia latente entre el bloque occidental del llamado \termino{Primer Mundo}, mayormente liderado por paises europeos y EUA con latinoamerica,
pues como lo describe \cite{topalov_cuestion}, el estado trabajó en conjunto con los actores sociales más importantes del momento en un esfuerzo por lograr una calidad de vida 
aceptable y digna, caso contrario en México, donde el estado si trabajó en conjunto con las instituciones existentes pero para el beneficio propio de la dictadura. \\
Por extraño que parezca, el control del partido hegemónico en el gobierno puso a un conjunto de personas que entre sus multiples fallas y altibajos, lograron consolidar una 
política pública que no se ha vuelto a repetir en el país, la creación de los multifamiliares a manos de Mario Panni son referentes del modelo urbanístico heredado por 
el urbanista Carlos Contreras. Estructuras como la de Tlatelolco o el multifamiliar Adolfo López Mateos fungieron como el reflejo de una modernidad simulada, pero un plan 
existente entre gobierno - sociedad. \\
A nivel transporte público, no se puede olvidar la interneción del ingeniero Bernardo Quintana Arrioja, quién a principios de 1960 crearía la llamada \termino{Solución Milán}. 
Que no sería otra cosa que un túnel del doble coraza que protegía las instalaciones en suelo fangoso, reforzando así cualquier medio de transporte profundo o muy profundo. 
Su implementación fue ignorada por el regente del entonces Distroto Federal Ernesto Peralta Uruchurtu, quién decidió que era un construcción costosa y estructuralmente 
peligrosa, además de inviable \cite{perlo2023} metropolis. Sin embargo fue hasta la designación de Alfonso Corona del Rosal en 1965 que se retomó el proyecto de 
un tren suburbano de alta capacidad, lo que se convertiría en el Metro de la Ciudad de México, al cual se le destinó un presupuesto federal para su posterior planeación. 
Y es aquí donde la entra la bifurcación del progreso, por un lado, la construcción de nuevas líneas durante todo el siglo XX ayudó a aliviar la creciente exploción demográfica 
que se estaba viviendo en esos años y por otro lado, se descuido la vivienda formal y se abandonaron los proyectos de vivienda a gran escala. \\
Tanto \cite{agustin_tragicomedia1} como \cite{molina2018} concuerdan en dos cosas; los movimientos sociales se sofocaron en la llamada \termino{Guerra Sucia} y el gobierno 
descuidó la planeación de vivienda en las periferias, creando lo que se conocería como \termino{Crecimiento Disperso} o \termino{Leap Frog} \cite{jia2010} 
por su nombre original en inglés. Ocasionando que aunque los ciudadanos de la creciente Ciudad de México pudiera moverse por el sistema de transporte colectivo Metro, 
sus viviendas estuvieran alejadas de un transporte masivo colectivo accesible y orillandolos a tiempos de traslado superiores a la hora y media para 1980 
\cite{salinas1980produccion}.

\subsection{El PRI en la Coyuntura}
\label{PRI-Coyuntura}

Durante las décadas de 1970 y 1980 existirían las primeras coyunturas políticas, socieles y económicas del PRI, sin embargo, antes de realizar un paseo por los sucesos más 
relevantes, no podemos empezar sin mencionar el análisis magnífico que hace \cite{bauman_estado}, quién define un periodo de \termino{Coyuntura} como una transición 
de una crísis política-social-económica a un nivel de estado pero con el teatro que siempre ha caracterizado a los gobiernos globalizados. Es decir, el arte de 
mencionar que existen pequeñas fisuras y periodos de transición y movimientos que debilitan la estructura; pero no la rompen. Claro que esto es solo discurso político.
Fue en la década de 1970 que el mundo transicionó de un modelo de estado dador de todo y administrador de todas las ramas de la economía a un simple intermediario para las 
multinacionales con poder económico similar a los PIB de países en desarrollo \cite{bauman_globalizacion_1999}.

\textbf{¿Y ahora qué?}\\
El problema explotó debido al desgaste y la mala administración a nivel global; como el mismo Bauman expresa; un solo hombre o una sola institución es incapaz de controlar
todos los ámbitos social-político-económico de una región por un periodo prolongado sin caer en la burocracía o la infeiciencia administrativa y México 
es el ejemplo perfecto de ello. Aunque el gobierno del presidente Luis Echeverría Álvarez fue polémico en muchas situaciones (represiones sociales), 
no se le puede negar que es el creador de muchas de las instituciones públicas que brindaron de dientes al urbanismo \cite{medin1982sexenio}: 
\begin{itemize} 
      \item \textbf{CONACYT:} (Consejo Nacional de Ciencia y Tecnología) - 1970: Creado para promover y articular las políticas públicas en materia de investigación 
      científica e innovación tecnológica.
      \item \textbf{INFONAVIT:} (Instituto del Fondo Nacional de la Vivienda para los Trabajadores) - 1972: Fundamental en la reconfiguración del espacio urbano moderno 
            en México, encargado de otorgar crédito para que los trabajadores adquieran vivienda, impulsando masivamente la expansión periférica de las ciudades.
      \item \textbf{UAM}: (Universidad Autónoma Metropolitana) - 1974: Fundada como un proyecto de desconcentración de la UNAM, con un modelo académico innovador y 
            unidades distribuidas estratégicamente (originalmente Azcapotzalco, Iztapalapa y Xochimilco) para articularse con distintas zonas de la ciudad.
      \item \textbf{FONACOT:} (Fondo Nacional para el Consumo de los Trabajadores) - 1974: Institución financiera para facilitar el acceso al crédito para bienes de 
            consumo a la clase trabajadora.
      \item \textbf{PROFECO:} (Procuraduría Federal del Consumidor) - 1976: Institución orientada a la defensa de los derechos de los consumidores, surgida tras la 
            promulgación de la Ley Federal de Protección al Consumidor.
\end{itemize}
Sin emabrgo, el sistema político y económmico estaba desgastado desde las entrañas, recordemos que el \termino{Milagro Mexicano} empezó con la instauración de la 
política económica de \textbf{Antonio Ortiz Mena} fue puesta en practica por primera vez en el sexenio de Adolfo Ruiz Cortinez en los años 50's, es decir; 
\textbf{habían pasado ya casi 20 años que el sistema económico y político mexicano no se actualizaba desde sus entrañas.} El estado de bienestar comenzaba 
a escacearle el dinero en el bolsillo y tuvo que cambiar la mira, no sin antes darse un tope de frente; la crísis de 1976 marcó el cierre definitivo, pues el presidente 
Echeverría se vió obligado a devaluar el peso al final de su mandato. \\
Tanto José Agustín como \cite{molina2018} concuerdan en que existieron dos personas que no eran las más aptas para el cargo de mandatario de estado: 
El ingeniero Pascual Ortiz Rubio que duró de 1930 a 1932, designado por el \termino{Jefe Máximo}; Plutarco Elías Calles y el licenciado José López Portillo o 
\termino{López por Pillo} \cite{agustin_tragicomedia1}. Una persona sin mucha experiencia política ni de estado designada por su antecesor. Y a pesar de ello, 
su gobierno al principio se enfrentó a una de las bonanzas económicas más inmensas de México en el siglo XX: El desubrimiento de los yacimientos de Canterel y 
Chicontepec. Los precios del petróleo y las inversiones públicas se elevaron por los cielos, pero la mala administración política y el rezagado sistema político y 
económico heredado no ayudaron en nada. \\
A nivel urbanístico esto retrasó dos factores clave:
\begin{itemize}
      \item \textbf{Social}: La gente perdió total fé en su gobierno, en sus gobernantes e impulsó la mira hacia otros proyectos de gobierno. Con el sistemA debilitado, 
            la mayor parte de las políticas públicas fueron de difícil implementación y una férrea resistencia que hasta al fecha no se ha podido reparar.
      \item \textbf{Económica:} El gobierno de Luis Echeverría Álvarez y José López Portillo marcaron el final de la bonanza mexicana, por lo que los proyectos ya hechos 
            con anterioridad, fueron desplazados a no importantes o intrascendentes y con el tiempo estos cayeron en desgracia o desuso; tal es el caso de los 
            multifamiliares de Mario Panni que durante buena parte de la década de 1980 presentaron un deterioro estructural junto con su posterior derrumbe a manos 
            del terremoto de 1985.
\end{itemize}
El tejido social se rompió y la consolidación Estado-Sociedad-Política inició el proceso de divorció y repartición de bienes.

\subsection{El PRI-MOR como resultado de la Coyuntura}

México como país no se puede entender sin la entrada de un nuevo sistema; El Neoliberalismo. Académicamente, el \textbf{neoliberalismo en México} 
se define como el proceso de reestructuración económica, política y social adoptado tras la crisis de la deuda externa de 1982, el cual marcó el fin del 
modelo de industrialización por sustitución de importaciones y del Estado benefactor interventor. Bajo la ortodoxia de los organismos financieros internacionales, 
este sistema se fundamentó en la reorientación de la economía hacia el libre mercado y la exportación. Sus ejes rectores incluyeron la desregulación financiera y comercial, 
la privatización masiva de empresas paraestatales, la contención salarial, la flexibilización laboral y una estricta disciplina fiscal orientada a reducir el déficit y el 
gasto público social \cite{tello2007estado}. \\
Anteriomente citamos al doctor en economía Carlos Salinas de Gortari en su obra \cite{salinas1980produccion}, este fue presentando como uno de los primeros informes en el 
gobierno de Miguel de la Madrid, donde él analizó y caracterizó uno de los problemas urbanos más latentes en el México de 1980 (y en la actualidad); la movilidad 
en la región. En este informe, se hizo una recolecta de los tiempos de traslado de los capitalinos de sus casas a sus trabajos y centro de estudios. Determinándose 
que el mexicano promedio se hacía de 1 hora a 3 horas de viaje en promedio, además de que el metro solo daba abasto al 30\% de la población de la Ciudad de México. 
Siendo que este no era ni el primero ni el segundo transporte de primera línea para los capitalinos (es decir el primero que abordas al salir de tu hogar) si no era el 
tercero. Esta consulta también sentó las bases para una de las politicas más sonadas en el siglo siguiente: \termino{Políticas de Género}. 
Pues en dicho documento se concentraba que la mayor parte de las personas que usaban el transporte público no eran del género masculino, quienes en su mayoría poseían 
automóvil particular. En cambio el sexo femenino representaba casi un 60\% de la carga del transporte y en su mayoría eran viajes concentrados en las siguientes actividades:
\begin{itemize}
      \item \textbf{Hogar - Centros Educativos}.
      \item \textbf{Hogar - Mercado o sucursal de abarrotes}. 
      \item \textbf{Hogar - Hogares colindantes}.
\end{itemize}
A partir de dicho informe, las políticas de género arrancaron en los años 90's orientados principalmente por un incremento en las llamadas \termino{Horas Hombre}. 
Sin embargo, México atravesaría por un periodo de esperanza, lo que \cite{molina2018} llama \termino{El especto Salinista} o \termino{El fantasma Salinista}. A nivel 
social, la percepción de la población mexicana crecía con gusto y las miradas internacionales se centraron en el \termino{TLC: Tratado de Libre Comercio}. 
Con el \termino{Salinismo}, México viró por completo al modelo Neoliberalista y se aferró a el como un ahogado a bolla de playa. \\
Durante los años 1988-1994, la estructura del Estado mexicano experimentó una profunda reingeniería institucional, alineada tanto con las políticas de liberalización 
económica (el inicio del neoliberalismo pleno) como con la necesidad de contener las presiones sociopolíticas surgidas tras los comicios de 1988. 
En este periodo se fundaron instituciones clave para la administración pública contemporánea, tales como:

\begin{itemize}
    \item \textbf{Instituto Federal Electoral (IFE, 1990):} Creado como respuesta a la crisis de legitimidad electoral, marcando el inicio de la ciudadanización en la 
            organización de los comicios (actualmente INE).
    \item \textbf{Comisión Nacional de los Derechos Humanos (CNDH, 1990):} Instituida originalmente como un órgano desconcentrado para la protección y defensa de los 
            derechos fundamentales frente a abusos del poder público.
    \item \textbf{Secretaría de Desarrollo Social (SEDESOL, 1992):} Creada en sustitución de la Secretaría de Programación y Presupuesto, concentrando las 
            atribuciones de política social, desarrollo regional, urbano y ecológico.
    \item \textbf{Instituciones Ambientales y Culturales:} Durante este periodo también se concretó la creación del Consejo Nacional para la Cultura y las Artes 
            (CONACULTA, 1988), la Procuraduría Federal de Protección al Ambiente (PROFEPA, 1992) y la Comisión Nacional para el Conocimiento y Uso de la 
            Biodiversidad (CONABIO, 1992).
\end{itemize}

En materia de política social, el régimen implementó el \textbf{Programa Nacional de Solidaridad (PRONASOL)}, el cual representó un cambio drástico de paradigma: 
el abandono del universalismo del Estado benefactor en favor de una política de \textit{focalización} \cite{barba2004paradigma}. 
Este enfoque buscaba mitigar la pobreza extrema generada por el ajuste estructural mediante transferencias e inversión comunitaria. Dicho modelo focalizado se convirtió 
en la columna vertebral de la política social en México durante casi tres décadas, conservando su esquema de transferencias condicionadas, aunque fue renombrado y 
modificado en sus reglas de operación por las administraciones subsecuentes \cite{gordon1999pobreza}:

\begin{enumerate}
    \item \textbf{PRONASOL} (Carlos Salinas de Gortari, 1988-1994).
    \item \textbf{PROGRESA} - \textit{Programa de Educación, Salud y Alimentación} (Ernesto Zedillo, 1997-2000).
    \item \textbf{Oportunidades} (Vicente Fox y Felipe Calderón, 2002-2012).
    \item \textbf{PROSPERA} (Enrique Peña Nieto, 2014-2018).
    \item \textbf{Programas para el Bienestar} (Andrés Manuel López Obrador, 2018-2024), donde se marcó un quiebre metodológico al transitar de las 
            transferencias \textit{condicionadas} a un esquema de apoyos directos de carácter universal por sectores demográficos.
\end{enumerate}
Sin embargo \termino{El Espectro Salinista} terminaría abruptamente con el \termino{Efecto Tequila}, proceso en el cual la economía de México se desplomó a niveles 
históricos, perdiendo por completo su presencia económica ganada durante principios de los años 90's. Este proceso lapidó la injerencia del partido hegemónico en 
el último pilar que aún controlaba: El régimen político. La ascención de Vicente Fox Quezada con el cambio de siglo enterró por completo las esperanzas de la perpetuidad. 
Pero por poético que suene esto, solo fue la evolución de un sistema que no ha acabado de cuajar; a partir del año 2000, la política mexicana ha quedado sin una 
representación fuerte y líderes políticos \termino{cojos}. Lo que se ha traducido en la implementación fallida de un nuevo tejido social: Lo urbano que necesita de 
una fuerte política que sea empujada desde altas esferas carece de apoyo de sus piernas económicas, las empresas no han podido consolidar un estado de equilibrio 
con los ciudadanos a quienes ven más como clientes y no como usuarios finales. La inserción laboral contemporánea ha evidenciado una falla sistémica estructural. 
Ante la ausencia de políticas públicas integrales que vinculen de manera efectiva al sistema educativo con el sector productivo y garanticen condiciones de trabajo 
dignas, los estudiantes y egresados se ven empujados hacia cuatro rutas críticas:

\begin{itemize}
    \item \textbf{La fuga de cerebros:} Ante la incapacidad del mercado interno para absorber mano de obra altamente calificada con salarios competitivos y 
            condiciones de desarrollo profesional, un sector importante de profesionistas opta por la emigración. Esta pérdida neta de capital humano acentúa 
            la dependencia tecnológica y limita la innovación científica del país \cite{delgado2011fuga}.
    
    \item \textbf{La deserción escolar:} Las apremiantes presiones económicas familiares, sumadas a la pérdida de confianza en la educación formal como vehículo 
            legítimo de movilidad social, propician el abandono prematuro de los estudios. Este fenómeno trunca el desarrollo de capacidades y perpetúa la 
            reproducción intergeneracional de la pobreza \cite{miranda2018desercion}.
    
    \item \textbf{El mercado informal y precario:} Una proporción mayoritaria de las juventudes que logran ingresar a la Población Económicamente Activa (PEA) 
            es absorbida por el sector informal. Esta vía se caracteriza por una flexibilización laboral extrema, la ausencia total de redes de seguridad social y 
            salarios de subsistencia, consolidando la precariedad como norma \cite{levy2008buenas}.
    
    \item \textbf{La incorporación a grupos criminales:} En contextos territoriales de alta marginación y vulnerabilidad institucional, la economía criminal se 
            erige como un empleador alternativo y un mecanismo ilusorio de estatus y movilidad rápida. Esta dinámica capitaliza la exclusión estructural de los 
            jóvenes, arrastrándolos a escenarios de violencia extrema y acortando drásticamente su esperanza de vida, un fenómeno conceptualizado desde la 
            sociología como \textit{juvenicidio} \cite{valenzuela2015juvenicidio}.
\end{itemize}
Con la llegada de Andrés Manuel López Obrador a la presidencia en 2018 a manos de su partido político \termino{MORENA: Movimiento de Regeneración Nacional}, 
la mira se puso en su aparato político; repetir un estado totalitarista en los tres poderes (Legislativo, ejecutivo y Judicial). Sin embargo, la historia no 
se repite, pero si rima.